\chapter{ComStock Building Models}
\label{chap:4_modeling}

ComStock uses about 30 high-level, whole-building characteristics to describe each building, as discussed in Section~\ref{chap:3_sampling}. However, whole-building energy models, such as the EnergyPlus\textsuperscript{\textregistered} model used by ComStock, typically require thousands of  inputs to describe a building for simulation. The purpose of the subsequent sections is to describe the assumptions, conventions, and data sources used to transform the high-level descriptions into inputs with the level of detail needed by EnergyPlus. Although the software used to implement this transformation is critical to the workflow, the focus is on the model inputs, not on the software workflow.

One question that often arises is why more of the input assumptions documented in this section are not incorporated directly into the sampling framework described in Section~\ref{chap:3_sampling}. This is an especially common question for those familiar with ResStock\textsuperscript{\texttrademark} \citep{resstock}, the residential building stock modeling tool that ComStock is based on. After all, ResStock uses more than 100 building characteristics to describe residential dwelling units, which are arguably less complex than commercial buildings. There are two main drivers behind the decision to limit the number of building characteristics: (1) handling complexity and (2) data availability for commercial buildings.

From a complexity standpoint, there is significantly more diversity among commercial buildings than among residential buildings. At one extreme, there are buildings like large hospitals, which may be several hundred thousand square feet, encompass spaces ranging from operating rooms to cafeterias, and be served by a complex array of HVAC systems. At the other extreme, there are buildings like small standalone retail stores, which may consist of just one retail space, a small storage room, and a restroom. Accounting for the diversity in lighting power density for each space type across all commercial building types in ComStock would alone require more than 100 building characteristics, many of which would not be applicable for certain building types. Multiply this by the number of characteristics that vary between building types, and the number of building characteristics required quickly becomes untenable.

From a data availability standpoint, there is simply much less information available for commercial buildings than there is for residential buildings. This means that modeling the commercial building stock requires more assumptions than modeling the residential building stock. Compounding this lack of data is the fact that most commercial building data sources handle complexity by focusing on a single building type (e.g., offices), providing information only at the whole-building level, or providing percentages of floor area associated with a given characteristic. Rather than making engineering estimates to generate probability distributions for every building characteristic, we have chosen to make point estimates for certain parameters. Proponents of stochastic modeling may disagree with this approach, but we believe it is warranted, given the model complexity that is avoided.

The end result is that many of the intra-building characteristics of commercial buildings must be inferred from whole-building characteristics. Rather than adding these to the input layer, they are set in the process that expands these whole-building characteristics into energy model inputs. That process is organized around a building specification, described next.

\section{The Building Specification}
\label{sec:building_spec}

Every ComStock model is generated from a single building specification that is assembled at run time and handed to ComStock-Typical (Section~\ref{sec:openstudio_standards}). The specification has two sources. The first is a building type definition, maintained with the ComStock source code, that records everything ComStock asserts about a building type that the space type data does not already carry. The second is the set of sampled characteristics for the individual model. The two are merged into one specification, which the library validates before it modifies the model.

A building type definition contains:
\begin{description}
\item[Space type ratios] The space types that compose the building and the fraction of floor area each occupies (Section~\ref{sec:space_type_ratios}). A building type may carry several entries, keyed by building subtype and by template family, so that, for example, a large office with a data center and one without, or a school built to a pre-2004 code and one built to a later code, can have different compositions. Each entry may also flag a space type as the \emph{default} (which absorbs floor area released when a very small space type is dropped), as the \emph{circulation} space type (which is placed as a double-loaded corridor), as having no geometry of its own, or as having its own story height or window-to-wall ratio.
\item[Form defaults] The aspect ratio, window-to-wall ratio, story height, and, where it differs, first-story height that apply when a sampled characteristic does not override them.
\item[Constructions] The exterior wall, roof, and floor construction types and any surface-level construction selections, by template family, that select assemblies from the standards construction data. A set of construction defaults shared by every building type is merged beneath the building type's own entries.
\item[Overrides] Building-type-specific departures from the space type data, grouped by family: thermostat overrides (for example, no unoccupied setbacks anywhere in a hospital), schedule overrides (a nurses station that keeps its own hours rather than the building's, a patient room with a lower lighting base), occupancy overrides (a hospital cafeteria denser than a restaurant dining room), service water heating overrides (a hospital kitchen that draws far less hot water per square foot than a restaurant kitchen), exhaust overrides, and load overrides. Each override names a space type, or the wildcard that matches every space type, and states only the fields it changes.
\end{description}

At run time, the sampled characteristics are expanded into the same specification: floor area, number of stories, aspect ratio, window-to-wall ratio, rotation, climate zone, weather file and design days, HVAC system type, weekday and weekend hours of operation, thermostat setpoints and setbacks, lighting and plug load base-to-peak ratios, wall construction type, interior lighting technology generation, service water heating fuel, and refrigeration technology level. Several of these enter as additional overrides rather than as post-processing of a finished model: the sampled setpoints and setbacks become thermostat overrides, and the sampled base-to-peak ratios become schedule overrides, both applied before any schedule or thermostat is built. A building that mixes building types---a strip mall with a restaurant, for example---is specified as a floor-area-weighted mix of two building type definitions, and the overrides of every building type in the mix are carried along.

Because the loads are keyed on the space type, most of what formerly distinguished one building type from another is now a consequence of its space type composition. Building type still determines the HVAC system type distributions (Section~\ref{sec:HVAC_System_Type}), the infiltration coefficients (Section~\ref{sec:infiltration_coefficients}), elevator counts, whether the service water heating loop recirculates, kitchen cooking equipment counts, refrigeration equipment, and the construction types named in the definition. Everything else---occupant density, lighting, plug and process loads, outdoor air, exhaust, water draws, thermostat setpoints, and schedules---comes from the space type, modified only by the overrides listed above.
